\honest{Incremental Progress and the Valley}{Eliezer Yudkowsky}{2009}

Yesterday I said that rationality is systematized winning. \nb{``Rationality is Systematized Winning'' is not in the book.} You object that reasonable people do not always win. Lottery winners are selective reporting; statistically you would never meet one. Rational agents can lose; they just cannot expect to do worse than some strategy they could have followed.

You mean something else: overconfident founders get rich, and religious people are happier. Here I concede. The optimality theorems hold for perfect reasoners, and no theorem says that every step toward the ideal improves performance. It ``is not, in fact, true.'' Then why strive, if we will never reach the ideal? A step backward has no guarantee either; judgment under uncertainty is what rationality is about. \nb{Economists had shown a close relative of this in 1956, the theory of the second best. The book itself had said in ``Knowing About Biases Can Hurt People'' (2007): ``If you're irrational to start with, having more knowledge can hurt you.''}

So why take the step? ``I can't necessarily answer for everyone.'' I give three reasons of my own. My work deals with deeply confused problems where one small mistake can lead you astray for years, so I must do better or go home. You may object that not everyone leads that kind of life. Refusing one step up forfeits all the steps beyond it, as Robin Hanson says that falling one step on the stairs leads to falling the next. Again you may object that not everyone is pushing the art into new territory. And once I have seen through a deception I cannot unsee it; I could no more believe in God than believe the sky green while looking at it. If you know enough to know you are better off deceiving yourself, it is too late. You may ask whether, knowing it may make my readers unhappier, I sponsor the collapse of their doublethink out of sadism.

My general answer is that my experience ``does seem to indicate'' that once you stop doing so many other things wrong, more rationality leaves you better off. In answering confusing questions, where I have practiced most, an optimist would be ``stomped into the ground'' by ``a casual student.'' I give no example. My essays reflect only that one area, not fighting akrasia, coordinating groups or being happy. For motivation and happiness I promise nothing. If those parts of the art have been developed systematically, I do not know of it. \nb{Eight days earlier, in ``Can Humanism Match Religion's Output?'', Yudkowsky had written that cognitive-behavioral therapy was ``experimentally shown to yield real improvements.''}

Even there I have tried to dispel half-mistaken ideas that get in a beginner's way: that rationality opposes feeling or value, or that sophisticated thinkers should be angsty and cynical. If someone develops those arts as far as I have developed the answering of impossible questions, I expect those wrapped in illusions will not begin to compete. I find it hard to believe that the strongest entrepreneur is wrapped in comforting overconfidence, or that the happiest life involves a tiny awareness that it is all a lie. I would sooner stake my hopes on neurofeedback or Zen meditation, though I have tried neither. \nb{Psychologists had disputed this since Taylor and Brown argued in 1988 that positive illusions foster mental health. A 2019 meta-analysis found self-enhancement linked to better adjustment in ordinary people.}

Two readers' comments show that the valley is real. One lost a girlfriend after deconverting and has been miserable and ineffective since, yet still prefers this to happiness with false beliefs. Another recovered from the same loss when a sister trained in CBT pointed out that it was not rational to expect to be alone forever. The first steps can hurt, and nothing guarantees you end up higher. But arguments against losing your illusions ``generally'' consider one step only, and ``the surveys'' compare the average religious person with the average atheist, not the most advanced theologians with the most advanced rationalists. \nb{No argument or survey is cited. Some studies that separate weak from strong belief suggest worse mental health in the middle than at either end, a dip like the post's valley.}

If you do not prefer truth to happiness with false beliefs, and do nothing precarious or confusing, buy no lottery tickets, and have already signed up for cryonics, then the steps are not guaranteed to leave you better off. If you do not care about the truth, have nothing to protect, do not want to push the art, and feel that your well-being depends on illusions, you are probably not reading this. If you are: sign up for cryonics, and stop reading before your illusions collapse: ``RUN AWAY!'' \nb{The post treats signing up for cryonics as the rational choice. The case for it is in ``You Only Live Twice,'' which the book does not include.}
